\setauthor{\firstauthor}
\subsection{Fragenkonfigurator}
Das Fragen-System der Matura-Vorbereitungsplattform ermöglicht den Nutzern, gezielt Fragen zu erstellen, zu verwalten und zum Lernen zu verwenden. Dabei unterscheiden sich die Rechte der Akteure deutlich:

\subsubsection{Schüler:innen-/User-Rechte}
\begin{itemize}
    \item Schüler:innen haben eingeschränkte Rechte innerhalb des Systems.
    \item Sie können eigene Fragen erstellen. Dabei haben sie zwei Möglichkeiten:
    \begin{itemize}
        \item \textbf{Private Fragen:} Diese bleiben nur für die/den jeweilige:n Schüler:in sichtbar und dienen zur persönlichen Vorbereitung.
        \item \textbf{Fragen mit Freigabeanfrage:} Schüler:innen können beim Erstellen einer Frage direkt eine Freigabe anfragen, damit diese von Lehkräften geprüft und möglicherweise öffentlich gemacht wird.
    \end{itemize}
    \item Schüler:innen können ihre eigenen privaten Fragen bearbeiten und löschen.
    \item Alle öffentlichen Fragen sind für Schüler:innen sichtbar, können jedoch nur von Lehkräften bearbeitet oder gelöscht werden.
    \item Beim Lernen können Schüler:innen zwischen öffentlichen Fragen und ihren eigenen privaten Fragen wählen.
    \item Der Lernmodus gibt keine Rückmeldung über richtige oder falsche Antworten, um ein stressfreies Wiederholen zu ermöglichen (in der Diplomarbeit \textit{Maturatrainer} ist die Überprüfung der Antworten im Lernmodus jedoch implementiert).
\end{itemize}

\subsubsection{Rechte der Lehkräfte}
\begin{itemize}
    \item Lehkräfte haben vollen Zugriff auf den öffentlichen Fragenpool.
    \item Sie können:
    \begin{itemize}
        \item neue öffentliche Fragen direkt erstellen,
        \item neue private Fragen für sich selbst erstellen,
        \item alle öffentlichen Fragen bearbeiten oder löschen,
        \item eigene private Fragen bearbeiten oder löschen,
        \item Fragen von Schüler:innen, für die eine Freigabe angefragt wurde, prüfen und öffentlich machen.
    \end{itemize}
    \item Lehkräfte haben Zugriff auf alle öffentlichen Fragen sowie auf Fragen, für die eine Freigabe angefragt wurde.
    \item Sie verwalten die Qualität und Sichtbarkeit der öffentlichen Inhalte im System.
\end{itemize}

\subsubsection{Funktionale Übersicht}
\begin{itemize}
    \item \textbf{Fragen erstellen:}
    \begin{itemize}
        \item Lehkräfte: können öffentliche Fragen direkt erstellen oder private Fragen für sich selbst.
        \item Schüler:innen: können private Fragen erstellen oder Fragen mit Freigabeanfrage erstellen.
    \end{itemize}

    \item \textbf{Fragen bearbeiten/löschen:}
    \begin{itemize}
        \item Lehkräfte: können alle öffentlichen Fragen sowie eigene private Fragen bearbeiten oder löschen.
        \item Schüler:innen: können nur ihre eigenen privaten Fragen bearbeiten oder löschen.
    \end{itemize}

    \item \textbf{Fragen freigeben:}
    \begin{itemize}
        \item Schüler:innen können beim Erstellen einer Frage eine Freigabe anfragen.
        \item Nur Lehkräfte können angefragte Fragen prüfen und bei öffentlich machen.
    \end{itemize}

    \item \textbf{Fragen anzeigen:}
    \begin{itemize}
        \item Öffentliche Fragen sind für alle Nutzer:innen sichtbar.
        \item Schüler:innen sehen zusätzlich ihre eigenen privaten Fragen.
        \item Lehkräfte sehen alle öffentlichen Fragen, Fragen mit Freigabeanfragen und eigene private Fragen.
    \end{itemize}

    \item \textbf{Fragen lernen:}
    \begin{itemize}
        \item Alle Nutzer:innen können Fragen im Lernmodus durchgehen.
        \item Schüler:innen können zwischen öffentlichen und ihren eigenen privaten Fragen wechseln.
        \item Der Lernprozess erfolgt ohne Rückmeldung über richtig oder falsch.
    \end{itemize}
\end{itemize}